% 竖版、横版样张 block-spacing 共用的正文，由两版目录的 block-spacing.tex 引入，不能单独编译。
% 检验代码卡片与引用块的前后距、二级列表与分隔线的竖向间距，断言见 preamble/tests/pdfcheck.py 的 check_block_spacing。
% 盒前的标记行（LINE-A1、LINE-A4）只用大写字母、连字符与平底数字（不用底部有弧的 3），
% 基线以下没有笔画，便于比较盒顶到基线的距离。
% 三组内容各从新的一页（双栏为新的一栏）开始，互不影响。
\section{Block spacing}
LINE-A1

\InputCode{bash}{preamble/tests/fixtures/sample.sh}
\InputCode{cpp}{preamble/tests/fixtures/示例 文件_1.cpp}
LINE-A2 卡片之后的正文。

LINE-A4
\begin{quoteblock}
QUOTE-A1 引用块。
\end{quoteblock}
\InputCode{bash}{preamble/tests/fixtures/sample.sh}
\begin{quoteblock}
QUOTE-A2 卡片之后的引用块。
\end{quoteblock}
LINE-A5 引用块之后的正文。
\newpage
\subsection{Headings}
\subsubsection{HEAD-B1}
\InputCode{bash}{preamble/tests/fixtures/sample.sh}
\subsubsection{HEAD-B2}
LINE-B1 卡片之后的标题，其后的正文。

LINE-B2
\subsubsection{HEAD-B3}
LINE-B3 正文之后的标题，其后的正文。
\newpage
LINE-C1
\begin{itemize}
  \item ITEM-C1 一级条目。
  \begin{itemize}
    \item ITEM-C2 二级条目。
    \item ITEM-C3 二级条目。
  \end{itemize}
  \item ITEM-C4 一级条目。
\end{itemize}
LINE-C2 列表之后的正文。

LINE-D1
\Separator{}
LINE-D2 分隔线之后的正文。
